\subsection{非営利組織の意思決定}

非営利組織、政府、公共部門では、利益最大化が目的ではない。目的は、一定の費用で社会的便益を最大化すること、または一定の目的を最小費用で達成することにある。

\subsubsection{便益費用分析}

便益費用分析は、提案がもたらす便益を費用で割り、比率で評価する方法である。便益費用比が 1 より小さい提案は、通常、費用の方が便益より大きいため退けられる。ただし、複数の提案を同時に考える場合には、追加的な分析が必要である。

\subsubsection{費用効果分析}

費用効果分析には二つの形がある。固定費用型では、決められた上限費用の中で便益を最大化する。固定効果型では、決められた成果を達成するための費用を最小化する。

\begin{notebox}{例}
自治体が住民向けオンライン申請システムを導入する場合、収益ではなく、待ち時間の削減、来庁回数の減少、職員負荷の軽減、住民満足度の向上などが便益になる。
\end{notebox}
